\section{Fallas, video largo y los límites de ingeniería en la sesión de Q\&A}

Después de las capacidades y del scaling, al final hay que volver a las fallas. Las colisiones complejas revelan que al modelo le falta una interacción estable entre varios cuerpos; el video largo está limitado por los tokens y por la attention cuadrática; el reasoning requiere traces y verifiers en forma de medios; el despliegue enfrenta además el watermark, la contaminación de datos y una serving architecture no publicada. Esta sección integra en la línea principal los límites expresados de forma oral en la sesión de Q\&A.

\subsection{Por qué las interacciones complejas todavía fallan}

La evaluación y el scaling muestran que el modelo es muy capaz en promedio; esta subsección se dirige a propósito a la parte menos promedio: las fallas en interacciones complejas. El ejemplo de colisión no es una anécdota graciosa, sino un diagnóstico: cuando varios objetos entran en contacto a gran velocidad, se ocultan entre sí y cambian sus relaciones topológicas, la plausibility visual local no basta para garantizar la coherencia del estado global. Al leer la figura, observe si las carrocerías se atraviesan, si el origen de los fragmentos es continuo y si la identidad de cada objeto se mantiene antes y después del choque.

\lecturefigure{slide-460-complex-motion-failure.jpg}{Caso de falla: en una colisión compleja de vehículos aparecen incoherencias geométricas y de interacción}{00:50:04}

\begin{warningbox}{Las fallas de video suelen ser fallas en la gestión del estado}
El modelo puede generar, cuadro por cuadro, texturas que «parecen una colisión» sin mantener un estado estable de los objetos, las restricciones de contacto ni el orden de los eventos. Aumentar la nitidez de cada cuadro no corrige este tipo de error por sí solo; se necesitan mejores representaciones espaciotemporales, datos, memoria de largo plazo, restricciones o mecanismos de autocorrección.
\end{warningbox}

\subsection{El video largo es, ante todo, un problema de longitud de secuencia}

Si se producen $n_f$ latent tokens por segundo y la duración es $S$, entonces $N\approx n_fS$. Con atención completa, el término dominante crece con la duración aproximadamente como
\[
\operatorname{cost}(S)=O((n_fS)^2d).
\]
Al pasar de 16 a 32 segundos, con las demás condiciones iguales, el número de tokens casi se duplica y el término dominante de attention se multiplica por cerca de cuatro; además, hay que mantener una trama, unos personajes y un estado de cámara más largos.

\begin{knowledgebox}{Los tres niveles de dificultad del video largo}
El primer nivel es el cómputo: crecen el context, las activaciones y la comunicación; el segundo son los datos: es más difícil obtener captions y límites de eventos para videos largos; el tercero es el modelado: el modelo debe mantener la identidad, las cadenas causales y el objetivo narrativo. Una memoria de GPU más grande solo alivia el primer nivel.
\end{knowledgebox}

\teachervoice{00:57:11--01:00:13, el docente dijo que 73K ya estaba cerca de la longitud de entrenamiento viable en ese momento. Si la duración se duplica, los costos de la secuencia y de la attention aumentan con rapidez; una película larga requiere además un diseño de sistema que vaya más allá de generar ingenuamente toda la secuencia de una sola vez.}

\subsection{Reasoning, verifier y native media generation}

La subsección anterior ya localizó los errores de gestión del estado en las interacciones complejas; esta se pregunta si el sistema puede detectarlos y corregirlos por iniciativa propia. El ponente propone dos direcciones abiertas: primero, que el modelo de video detecte errores evidentes y se autocorrija antes de entregar la salida; segundo, incorporar la generación de video en un LLM multimodal nativo. Ambas deben responder si el objetivo de aprendizaje es unificado y cómo verificar que una salida multimedia es «correcta», algo más difícil de definir que aumentar los parámetros del modelo.

\lecturefigure{slide-464-future-directions.jpg}{Direcciones futuras: continuar con el scaling, reasoning sobre medios y generación multimodal nativa}{00:53:37}

Si se escribe el ciclo generar--revisar--corregir, puede abstraerse como
\[
y^{(k)}=G_\theta(z,c,r^{(k)}),\qquad
e^{(k)}=V_\phi(y^{(k)},c),\qquad
r^{(k+1)}=R(r^{(k)},e^{(k)}),
\]
donde $G$ es el generador, $V$ es el verifier y $R$ actualiza el reasoning state. La dificultad es que el video no cuenta con un criterio de corrección barato y determinista como el de un problema de matemáticas; que una persona note de inmediato que un choque es inverosímil no significa que exista una reward automática entrenable.

\begin{importantbox}{El cuello de botella central del reasoning sobre medios es la verificación}
Generar «texto de razonamiento» no mejora el video automáticamente. Hay que definir errores detectables (identidad de los objetos, orden de las acciones, geometría, física, restricciones textuales o estética) y demostrar que el verifier se correlaciona con el juicio humano. De lo contrario, el reasoning trace puede ser solo tokens adicionales y no una autocorrección efectiva.
\end{importantbox}

\subsection{Cinco precisiones de la sesión de Q\&A}

La sesión final de preguntas y respuestas aclaró mejor los límites de la arquitectura, los datos y la gobernanza. Los cinco puntos siguientes no son un apéndice, sino condiciones necesarias para entender las conclusiones de la exposición principal.

\begin{enumerate}
  \item \textbf{Inductive bias:} las CNN pueden tener ventaja a pequeña escala o en ciertas diffusion policies para robots; el Transformer escala con más facilidad sobre una infraestructura madura, pero el debate no ha terminado.
  \item \textbf{Physics:} se pueden codificar priors de forma explícita, añadir datos de simulación o confiar en la escala; cuál ruta es la mejor sigue siendo una pregunta abierta, y los ejemplos actuales no demuestran que la física esté resuelta.
  \item \textbf{Synthetic data:} el contenido generado de baja calidad que entra en el preentrenamiento contamina los datos, pero los datos sintéticos no son dañinos por naturaleza; en el post-training controlado ya es habitual usar muestras generadas por modelos.
  \item \textbf{Safety:} los deepfakes y el watermarking son problemas independientes de investigación, desarrollo y despliegue, y no pueden sustituirse con métricas de calidad de generación.
  \item \textbf{Public evidence:} el artículo publica detalles de la infra de entrenamiento, pero en la clase no se presentó la serving architecture completa; no conviene especular cuántas GPU usa una sola solicitud.
\end{enumerate}

\begin{knowledgebox}{¿En qué nivel debe ubicarse la «unificación»?}
En la sesión de Q\&A, el docente señaló que, siempre que una nueva modalidad pueda codificarse en tokens, el núcleo Transformer puede reutilizarse; la dificultad real suele estar en la representación y en el objetivo. Un tronco unificado no exige unificar el tokenizer, el decoder, la loss ni el protocolo de evaluación.
\end{knowledgebox}

\teachervoice{00:54:10--00:56:03, el docente reconoce que el sesgo inductivo de las CNN especializadas puede ser valioso a pequeña escala; la experiencia de Movie Gen es que el Transformer a gran escala escala con más facilidad y cuenta con una infraestructura más madura, pero aclaró de forma explícita que no da por cerrado el debate.}

\teachervoice{01:09:23--01:10:33, sobre la contaminación por datos sintéticos, el docente distingue entre los videos falsos de baja calidad en el preentrenamiento y los datos controlados del post-training: los datos generados no son malos por naturaleza; lo que importa es la calidad, el objetivo y el filtrado.}

\subsection{Resumen de la sección}

Las fallas en interacciones complejas revelan que los modelos de video gestionan el estado de forma insuficiente; el video largo está limitado a la vez por el cómputo cuadrático, los datos y la coherencia de largo plazo; lo decisivo en el reasoning es el verifier, no el texto adicional. La sesión de Q\&A precisó además los límites del Transformer, de los priors físicos, de los datos sintéticos, de la seguridad y de la evidencia pública sobre el sistema.
